\documentclass[11pt]{article}
\usepackage[margin=1in]{geometry}
\usepackage{booktabs}
\usepackage{hyperref}
\usepackage{amsmath}

\title{A Criticality Claim That Does Not Survive Distribution Fitting:\\
the Empedoclean Node Field Re-examined}
\author{Giovanni Pezzin}
\date{}

\begin{document}
\maketitle

\begin{abstract}
An earlier internal report claimed that a deterministic node field with
coupled attraction and repulsion exhibits self-organised criticality: a
band of the repulsion parameter in which the cluster-size distribution
follows a power law. That claim rested on a single seed, one node count and
one time horizon. We re-run the sweep over 16 seeds and nine parameter
settings and apply standard distribution fitting: discrete power-law
maximum likelihood with the lower cutoff chosen by minimising the
Kolmogorov--Smirnov distance, a lognormal fit to the same tail, and the
Vuong log-likelihood-ratio test. The power law never wins. At three of the
seven fittable settings the lognormal is significantly better; at the
setting originally reported as critical the two are statistically
indistinguishable, which is the absence of evidence for criticality rather
than evidence for it. We further find that the regime label at the claimed
critical band is not stable across seeds. The original positive result was
a single-seed artifact of a heavy-tailed snapshot. We report the negative
result, keep the falsified claim visible, and state what a criticality
claim in a toy dynamical system would actually require.
\end{abstract}

\section{What was claimed}

The Empedoclean dyad is a deterministic dynamical system over a field of
nodes with two forces: an attractive term (``Eros'', strength
\texttt{eros\_strength}) and a repulsive term (``Nemesis'',
\texttt{nemesis\_strength}). The original report described a 300-node,
2-dimensional, 400-tick sweep of the repulsion strength against a fixed
attraction of 0.10, and reported three regimes: a supercritical regime in
which one giant cluster forms, a subcritical regime in which the field
fragments into singletons, and --- between them --- a \emph{critical} band
at $\text{nemesis} \approx 0.01$ exhibiting a power-law cluster-size
distribution. The reported evidence for that band was a single
cluster-size snapshot, whose ten largest clusters were
$59, 32, 9, 9, 8, 8, 7, 6, 6, 5$.

That evidence was always weak, and the original report said so:
``a heavy-tailed snapshot under one configuration, not an SOC proof.''
The open questions it listed were multiple seeds, finite-size analysis,
and distribution fitting with goodness-of-fit and confidence intervals.
This paper answers them.

\section{Method}

\subsection{Sweep}

We use the same field, the same geometry, and the same link distance as the
original: 300 nodes, 2 dimensions, 400 ticks, \texttt{eros\_radius} 0.30,
\texttt{eros\_strength} 0.10, \texttt{nemesis\_radius} 0.20,
\texttt{max\_step} 0.03, decay, death and synthesis disabled to isolate the
attraction/repulsion balance. We sweep
$\text{nemesis} \in \{0, 0.001, 0.003, 0.006, 0.01, 0.02, 0.04, 0.08, 0.16\}$
and run \textbf{16 seeds} per setting, yielding 144 independent runs.

\subsection{Fitting}

For each setting we pool cluster sizes across seeds and follow the
procedure of Clauset, Shalizi and Newman (2009):

\begin{enumerate}
\item Choose the lower cutoff $x_{\min}$ by minimising the
  Kolmogorov--Smirnov distance between the empirical tail and the fitted
  discrete power law.
\item Fit the discrete power-law exponent $\alpha$ by maximum likelihood,
  normalising with the Hurwitz zeta function $\zeta(\alpha, x_{\min})$.
\item Fit a lognormal to the same tail by maximum likelihood.
\item Compare the two non-nested models with the Vuong
  log-likelihood-ratio test. A positive $R$ favours the power law; a
  negative $R$ favours the lognormal; a non-significant $p$ means the two
  are indistinguishable on this data.
\end{enumerate}

This is the test the original report named as the condition under which
the criticality claim would be unfounded.

\section{Results}

\begin{table}[h]
\centering
\small
\begin{tabular}{rrrrrrrl}
\toprule
\textbf{nemesis} & \textbf{clusters} & $x_{\min}$ & $\alpha$ & \textbf{KS} & $R$ & $p$ & \textbf{verdict} \\
\midrule
0.000 & 50   & 76 & 3.558 & 0.091 & 4.2   & 0.291 & indistinguishable \\
0.006 & 81   & 2  & 1.201 & 0.534 & $-17.4$ & 0.046 & \textbf{lognormal better} \\
0.010 & 2507 & 82 & 4.872 & 0.166 & $-2.5$ & 0.396 & indistinguishable \\
0.020 & 4576 & 3  & 5.000 & 0.900 & $-0.9$ & 0.856 & indistinguishable \\
0.040 & 4722 & 2  & 5.000 & 0.987 & $-89.1$ & 0.014 & \textbf{lognormal better} \\
0.080 & 4747 & 2  & 5.000 & 0.918 & $-18.9$ & 0.013 & \textbf{lognormal better} \\
0.160 & 4758 & 2  & 5.000 & 1.000 & --- & --- & undetermined \\
\bottomrule
\end{tabular}
\caption{Distribution fits, 16 seeds per setting. Settings 0.001 and 0.003
produce too few clusters to fit. $R$ is the Vuong statistic; positive
favours the power law.}
\label{tab:fits}
\end{table}

\textbf{The power law never wins.} In no setting does $R$ come out
positive and significant. At three settings the lognormal is significantly
better than the power law. At the setting originally reported as critical
($\text{nemesis} = 0.010$), the two models are indistinguishable
($R = -2.5$, $p = 0.40$) --- which means this data provides no evidence
distinguishing a power law from a lognormal, not that a power law holds.

Two further observations bear on the fits. First, $\alpha$ runs to the
upper bound of the search in the fragmented regime and the KS distance
approaches 1.0, indicating that a power law is simply the wrong family
there. Second, $x_{\min}$ chosen at 2 or 3 in the fragmented settings means
the fitted tail consists of almost all clusters, including the abundance of
singletons --- the classic failure mode where a heavy-tailed-looking
distribution is fitted above a cutoff so low that the fit is meaningless.

\subsection{The critical band is not stable across seeds}

The original claim located a critical band at $\text{nemesis} \approx 0.01$.
Classifying each of the 16 runs at that setting with the code's own regime
rule yields a mixture: some runs classify as \emph{Critical} and some as
\emph{Subcritical}. At no other setting does the classifier disagree across
seeds. The mean largest-cluster fraction moves from 0.986 at
$\text{nemesis} = 0.006$ to 0.346 at 0.010 to 0.009 at 0.020, so what the
sweep actually shows is a \emph{sharp} transition between two phases, not a
self-organised band in which the system sits at a critical point. A sharp
transition whose location varies with the seed is the opposite of the
robustness that self-organised criticality is supposed to explain.

\section{Why the original positive was an artifact}

The original snapshot had ten clusters, of which the largest was 59 and the
rest fell quickly. A snapshot of that shape is exactly what a lognormal
with a heavy tail produces, and a handful of cluster sizes cannot
distinguish the two families. With one seed there is also no way to see
that the regime label is unstable, and with $n = 10$ clusters there is no
power to fit a tail at all. The artifact was not a coding error: the
measurement was correctly taken, and correctly flagged as insufficient.
It became a claim only in the retelling.

\section{What a criticality claim would require}

We state this so the refutation has a constructive form. A defensible claim
of criticality in a system like this needs, at minimum:

\begin{enumerate}
\item \textbf{Many seeds}, with the regime label stable across them at the
  claimed setting.
\item \textbf{Distribution fitting} against at least a lognormal
  alternative, with the cutoff chosen by a principled criterion rather than
  by eye, and goodness-of-fit reported.
\item \textbf{A finite-size analysis}: the cutoff behaviour of the
  cluster-size distribution should change with node count in a specified
  way if the system is critical.
\item \textbf{An order parameter and its scaling}, not merely a
  heavy-tailed snapshot.
\item \textbf{Control systems} whose exact critical point is known, run
  through the same analysis pipeline, to show the pipeline can detect
  criticality when it is present.
\end{enumerate}

Point 5 is the one we regard as most valuable and least often done: without
a positive control, a null result cannot be distinguished from an
insensitive instrument. We have not run it, so we cannot claim our pipeline
would have detected a genuine power law; we claim only that it does not
detect one here.

\section{What survives}

The system is not worthless because the criticality claim failed. What
survives is narrower and better supported:

\begin{itemize}
\item The attraction/repulsion balance produces a sharp transition between
  a collapsed phase and a fragmented phase, stable across 16 seeds at every
  setting except the transition region.
\item The transition region is seed-sensitive, which is a property worth
  knowing before anyone builds on the system at that setting.
\item The promotion gate the project had already imposed on this module
  (gate 2, no-regression, still unpassed) was the right call; the module
  remains research-only and touches no production path.
\end{itemize}

\section{Reproducibility}

\begin{itemize}
\item Sweep: \texttt{examples/empedocles\_multiseed.rs}, run as
  \texttt{cargo run --release --features empedocles --example
  empedocles\_multiseed -- 16}, writing
  \texttt{research/empedocles\_multiseed.json}.
\item Fits: \texttt{research/papers/fit\_empedocles.py}, writing
  \texttt{research/empedocles\_fits.json}.
\item The example is declared in \texttt{Cargo.toml} with
  \texttt{required-features = ["empedocles"]}, matching the two pre-existing
  Empedocles examples.
\item The original single-seed example,
  \texttt{examples/empedocles\_criticality.rs}, is unchanged and still
  builds. Its output is the artifact this paper refutes; we have kept it
  rather than deleted it.
\end{itemize}

\section{Conclusion}

A single-seed heavy-tailed snapshot was reported as evidence of
self-organised criticality in a 300-node dynamical system. Re-running the
sweep with 16 seeds and fitting the resulting distributions properly shows
that the power law never beats a lognormal, that at the claimed critical
setting the two are indistinguishable, and that the regime label at that
setting is unstable across seeds. The claim is refuted. The system's
interesting behaviour --- a sharp, seed-influenced transition between
collapsed and fragmented phases --- survives, but it is a phase transition,
not self-organised criticality. We keep the refuted claim, the original
example and the original report in the repository, because the sequence
from plausible snapshot to refuted claim is the most useful thing this line
of work produced.

\begin{thebibliography}{9}
\bibitem[Clauset et al.(2009)]{clauset2009powerlaw}
Clauset, A., Shalizi, C.~R., \& Newman, M.~E.~J. (2009).
\newblock Power-law distributions in empirical data.
\newblock \emph{SIAM Review}, 51(4), 661--703.
\newblock DOI: 10.1137/070710111.
\end{thebibliography}

\end{document}
